\chapter{Zadatak 2 --- priprema za obranu}
\label{ch:z2-obrana}

\section{Brojevi koje moraš znati napamet}
\label{sec:z2o-brojevi}

\begin{longtable}{@{}p{0.32\linewidth} p{0.30\linewidth} p{0.32\linewidth}@{}}
\toprule
\textbf{Veličina} & \textbf{Egzaktno} & \textbf{Decimalno} \\
\midrule
\endhead
Intervali koeficijenata & --- & $[-1{,}15,\,-0{,}85]$, osim $a_{21} \in [0{,}85,\,1{,}15]$ \\
Broj vršnih nejednakosti & $2 \cdot 2^2$ & $8$ \\
Najgore prvo ograničenje & $x_1 + x_2 \geq 200/17$ & $11{,}764706$ \\
Najgore drugo ograničenje & $1{,}15x_1 - 0{,}85x_2 \leq 3$ & --- \\
$x_{\mathrm{rob}}^{\star}$ & $(13/2,\ 179/34)$ & $(6{,}500000,\ 5{,}264706)$ \\
$f(x_{\mathrm{rob}}^{\star})$ & $978253/11560$ & $84{,}623962$ \\
$\grad f(x_{\mathrm{rob}}^{\star})$ & --- & $(13{,}420588,\ 16{,}108824)$ \\
Aktivni vrhovi & $k = 1$ i $k = 7$ & $(-0{,}85,-0{,}85)$, $(1{,}15,-0{,}85)$ \\
Množitelji & --- & $\mu_1^{\star} = 17{,}607439$, $\mu_7^{\star} = 1{,}344118$ \\
Odbačeni kandidat & --- & $(6{,}9073,\ 4{,}8574)$, $g_2^{\mathrm{wc}} = +0{,}8146$ \\
Cijena robusnosti & --- & $+24{,}017901$ $(+39{,}63\%)$ \\
Monte Carlo kršenja & --- & nominalno $54/100$, robusno $0/100$ \\
\bottomrule
\end{longtable}

\begin{izgovor}[Trik za pamćenje]
$x_1 = 6{,}5$ je \emph{isti broj} kao prva komponenta vrha dopustivog skupa u
Zadatku 1. Nije slučajnost --- oba puta dolazi iz istog razloga: zbrajanjem
dviju aktivnih jednadžbi članovi s $x_2$ ispadnu jer imaju suprotne predznake.
U Zadatku 1 zbroj koeficijenata uz $x_1$ je $1 + 1 = 2$ i desne strane
$10 + 3 = 13$; ovdje je $0{,}85 + 1{,}15 = 2$ i opet $10 + 3 = 13$. Skaliranje
za $\pm15\%$ simetrično je oko jedinice, pa se zbroj ne mijenja.
\end{izgovor}

\section{Zadatak u četiri rečenice}
\label{sec:z2o-sazetak}

\begin{izgovor}[Ako imaš samo minutu]
,,Koeficijenti ograničenja nisu poznati točno, nego samo unutar $\pm15\%$ oko
nominalnih. Skup nesigurnosti za svaki redak je pravokutnik, jer koeficijenti
variraju nezavisno.

Uvjet mora vrijediti za svaku realizaciju, što je beskonačno mnogo nejednakosti.
Sažeo sam ih u uvjet na supremum, a onda iskoristio da linearna funkcija po
konveksnoj ljusci postiže maksimum u vrhu. Time robusno ograničenje postaje
egzaktno ekvivalentno s četiri nejednakosti po ograničenju, ukupno osam.

Robusni problem je time običan kvadratni program s osam ograničenja. Aktivna su
dva --- najgori vrhovi --- pa je rješenje šest i pol, pet zarez dva šest pet, s
oba množitelja strogo pozitivna.

Cijena robusnosti je porast funkcije cilja za trideset devet i pol posto. U
uzorku od sto realizacija nominalno rješenje krši ograničenja u pedeset četiri
slučaja, robusno ni u jednom.``
\end{izgovor}

\section{Pitanja koja se očekuju}
\label{sec:z2o-pitanja}

\begin{pitanje}{Zašto je dovoljno provjeriti samo vrhove?}
\textbf{Odgovor.} Jer je skup nesigurnosti pravokutnik, dakle konveksna ljuska
svojih vrhova, a funkcija $a \mapsto a\trans x$ je linearna u $a$. Svaku točku
skupa zapišem kao konveksnu kombinaciju vrhova; uvrštavanjem i korištenjem
nenegativnosti težina dobiva se da je vrijednost u toj točki najviše jednaka
najvećoj vrijednosti u vrhovima. Obrat je trivijalan jer su vrhovi elementi
skupa. Ekvivalencija je egzaktna, nije aproksimacija.
\end{pitanje}

\begin{pitanje}{Bi li to vrijedilo da je skup nesigurnosti elipsoid?}
\textbf{Odgovor.} Ne bi --- elipsoid nema vrhova. Supremum se tada računa
analitički: za $\mathcal{P} = \{a_0 + Pu \mid \norm{u}_2 \leq 1\}$ vrijedi
$\sup_a a\trans x = a_0\trans x + \norm{P\trans x}_2$, pa robusno ograničenje
postaje $a_0\trans x + \norm{P\trans x}_2 \leq h$. To više nije linearno nego
konusno ograničenje drugog reda, i problem postaje SOCP. Predavanje 8 to obrađuje
na zadnja dva slajda.
\end{pitanje}

\begin{pitanje}{Zašto je za prvo ograničenje najgori koeficijent $-0{,}85$, a ne $-1{,}15$?}
\textbf{Odgovor.} Jer je ograničenje oblika $a\trans x \leq -10$, dakle traži da
lijeva strana bude \emph{dovoljno negativna}. Uz $x \geq 0$, koeficijent bliži
nuli čini lijevu stranu manje negativnom, dakle uvjet teže zadovoljivim. Najgori
je stoga $-0{,}85$. Za drugo ograničenje, koje traži da lijeva strana ne
premaši $3$, najgori je najveći koeficijent uz $x_1$, dakle $1{,}15$. Pravilo je
uvijek isto: bira se ono što \emph{maksimizira lijevu stranu}.
\end{pitanje}

\begin{pitanje}{Zašto niste nabrojali svih $2^8 = 256$ kombinacija aktivnosti, kao u Zadatku 1?}
\textbf{Odgovor.} Jer se unaprijed vidi da šest od osam nejednakosti ne mogu biti
vezujuće. Uz pretpostavku $x \geq 0$, za svako ograničenje jedan vrh implicira
ostala tri, pa preostaju samo dva kandidata. Provjerio sam obje mogućnosti:
samo prvi aktivan daje nedopustivu točku, pa moraju biti aktivna oba. Na kraju
sam provjerio i sve preostale nejednakosti u tablici --- sve imaju strogu
rezervu. Pretpostavka $x \geq 0$ potvrđena je dobivenim rješenjem.
\end{pitanje}

\begin{pitanje}{Koliko bi vrhova bilo da nesigurnost ima dvadeset koeficijenata?}
\textbf{Odgovor.} $2^{20}$, dakle preko milijun nejednakosti po ograničenju ---
metoda vrhova bi postala neupotrebljiva. Tada se koristi dualna reformulacija s
Predavanja 8, slajd 3/21, koja dodaje samo onoliko novih varijabli koliko skup
nesigurnosti ima nejednakosti u svom opisu, dakle linearno mnogo. Obje su metode
egzaktne; razlikuju se samo po veličini rezultirajućeg problema.
\end{pitanje}

\begin{pitanje}{Je li robusno rješenje ,,bolje`` od nominalnog?}
\textbf{Odgovor.} Ne po vrijednosti funkcije cilja --- ondje je gore za
$39{,}6\%$, i to nužno, jer je robusni dopustivi skup podskup nominalnoga. Bolje
je po pouzdanosti: dopustivo je za svaku realizaciju, dok je nominalno dopustivo
za otprilike polovicu. Koji je izbor ispravan ovisi o tome što se plaća skuplje:
$24$ jedinice funkcije cilja ili $50\%$ vjerojatnost da ograničenje bude
prekršeno.
\end{pitanje}

\begin{pitanje}{Što bi se dogodilo da je nesigurnost $\pm 30\%$?}
\textbf{Odgovor.} Omotnice bi se dodatno pomaknule --- najgore prvo ograničenje
postalo bi $x_1 + x_2 \geq 10/0{,}70 = 14{,}286$ --- pa bi robusni optimum bio
još dalje, a cijena robusnosti veća. Struktura rješenja ostala bi ista: opet
osam vršnih nejednakosti, opet dva aktivna vrha. Pri dovoljno velikoj
nesigurnosti robusni problem može postati i nedopustiv, ali ovdje ne, jer je
dopustivi skup neograničen.
\end{pitanje}

\section{Popis zamki}
\label{sec:z2o-zamke}

\begin{longtable}{@{}p{0.05\linewidth} p{0.42\linewidth} p{0.47\linewidth}@{}}
\toprule
\textbf{\#} & \textbf{Zamka} & \textbf{Kako je izbjeći} \\
\midrule
\endhead
1 & Interval zapisan u krivom poretku za negativnu nominalnu vrijednost &
$a_{11} \in [-1{,}15,\,-0{,}85]$, ne obrnuto (\S\ref{sec:z2-intervali}). \\[5pt]
2 & Brkanje skupa nesigurnosti i dopustivog skupa & Prvi je u prostoru
koeficijenata, drugi u prostoru varijabli. \\[5pt]
3 & Jedan zajednički skup nesigurnosti za obje matrice retke & Skupovi su
neovisni, $\mathcal{P}_1$ i $\mathcal{P}_2$ zasebno. \\[5pt]
4 & ,,Najgori je najveći koeficijent`` & Najgori je onaj koji maksimizira lijevu
stranu; ovisi o predznaku. \\[5pt]
5 & Vrhovi obrazloženi intuicijom umjesto dokazom & Argument mora ići preko
konveksne kombinacije (\S\ref{sec:z2-teorem-vrhovi}). \\[5pt]
6 & Pretpostavka $x \geq 0$ koja se ne provjeri & Rješenje ima obje komponente
pozitivne --- to treba izgovoriti. \\[5pt]
7 & Zaboravljeni preostali vrhovi & Svih osam nejednakosti treba provjeriti;
rezerve su u tablici. \\[5pt]
8 & $0/100$ predstavljeno kao dokaz & Dokaz su vršne nejednakosti; realizacije
ilustriraju. \\[5pt]
9 & Uniformna razdioba predstavljena kao dio zadatka & Zadatak ne spominje
razdiobu; ona je odabir za crtanje. \\[5pt]
10 & Tvrdnja da je robusno rješenje ,,bolje`` & Bolje po pouzdanosti, lošije po
funkciji cilja --- i to nužno. \\
\bottomrule
\end{longtable}

\section{Veza sa Zadatkom 1}
\label{sec:z2o-veza}

\begin{predznanje}
Zadatak 2 je Zadatak 1 s više ograničenja. Sve što se prenosi:
\begin{itemize}[nosep,leftmargin=*]
  \item ista funkcija cilja, ista matrica $Q$, isti $c$ --- pa iste nivo-elipse i
        isto središte $x^{0}$;
  \item ista provjera konveksnosti: $Q \pd 0$, ograničenja afina;
  \item isti KKT aparat, samo s osam umjesto dva ograničenja;
  \item ista logika aktivnog skupa: aktivna ograničenja imaju $\mu > 0$,
        neaktivna $\mu = 0$;
  \item ista interpretacija množitelja --- $\mu_1^{\star} = 17{,}61$ znači da bi
        popuštanje najgore verzije prvog ograničenja za jedinicu smanjilo
        optimalnu vrijednost za otprilike $17{,}61$.
\end{itemize}
\end{predznanje}
